\documentclass{scrartcl}
\usepackage[utf8]{inputenc}
\usepackage[T1]{fontenc}
\usepackage{amsmath}
\usepackage{hyperref}
\usepackage{url}
\usepackage{natbib}
\usepackage{graphicx}
\usepackage{listings}
\usepackage{booktabs}
\usepackage{xcolor}
\usepackage{enumitem}
\usepackage{tikz}
\usetikzlibrary{positioning, arrows.meta, fit, backgrounds, calc,
                shapes.geometric, shapes.multipart}
\usepackage{float}
\usepackage{cleveref} % must be last

\definecolor{lstbg}{rgb}{0.97,0.97,0.97}
\definecolor{lstkw}{rgb}{0.00,0.40,0.70}
\definecolor{lstcm}{rgb}{0.35,0.55,0.35}
\definecolor{lststr}{rgb}{0.65,0.15,0.15}
\definecolor{lstln}{gray}{0.55}

\lstdefinestyle{scopa}{
    backgroundcolor=\color{lstbg},
    basicstyle=\ttfamily\footnotesize,
    keywordstyle=\color{lstkw}\bfseries,
    commentstyle=\color{lstcm}\itshape,
    stringstyle=\color{lststr},
    showstringspaces=false,
    breaklines=true,
    breakatwhitespace=true,
    numbers=left,
    numberstyle=\tiny\color{lstln},
    numbersep=6pt,
    tabsize=2,
    captionpos=b,
    frame=single,
    framerule=0.3pt,
    rulecolor=\color{gray!50},
}
\lstset{style=scopa}

\newcommand{\emailaddr}[1]{\href{mailto:#1}{\texttt{#1}}}

% Compilable figure placeholder: framed box with descriptive caption.
% Usage: \figph{label}{height-in-cm}{caption text}
\newcommand{\figph}[3]{%
    \begin{figure}
        \centering
        \fbox{\begin{minipage}{0.82\textwidth}\centering\vspace{#2}\textbf{Figure placeholder:} \texttt{#1.png}\\[0.4em]\textit{To be drawn from the description below.}\vspace{#2}\end{minipage}}
        \caption{#3}
        \label{fig:#1}
    \end{figure}%
}

% TikZ figure include: pulls diagrams/<label>.tex, scales to text width, captions, labels.
% Usage: \figtikz{label}{caption text}
\newcommand{\figtikz}[2]{%
    \begin{figure}
        \centering
        \resizebox{\textwidth}{!}{\input{diagrams/#1.tex}}
        \caption{#2}
        \label{fig:#1}
    \end{figure}%
}

\title{\LARGE
    Scopa2\\
    A Geo-Distributed Multiplayer Card Game
}
\subtitle{Final Report for the Distributed Systems Course --- AY 2025/2026}

\author{
    Marco Coppola \\ \emailaddr{marco.coppola3@studio.unibo.it}
    \and
    Valerio Pio De Nicola \\ \emailaddr{valeriopio.denicola@studio.unibo.it}
}

\date{\today}

\begin{document}

\maketitle

\begin{abstract}
\emph{Scopa2} is a multiplayer turn-based card game built as a vehicle for
exercising, end-to-end, the engineering principles of distributed systems.
The system pairs two players from anywhere in the world through an
ELO-bracketed matchmaking queue, plays out a custom variant of the Italian
card game \emph{Scopa Scientifica} extended with a card-mutation mechanic
inspired by liturgical \emph{Santi}, and broadcasts every move in real time
across a WebSocket fabric backed by a Pusher-compatible message bus. Two
geographic regions (Europe and North America) run a fully symmetric stack
on Kubernetes; clients reach the closest one through a latency-aware
discovery routine and survive both single-node and whole-region failures
through a combination of stateless application replicas, a Sentinel-managed
Redis backplane, and a distributed YugabyteDB cluster with intra-region
Raft consensus and asynchronous cross-region replication. The report
discusses the design choices in light of the CAP/PACELC trade-off,
delegates consensus to layers that have already solved it, and turns the
game's deterministic state machine into a recovery mechanism via
event-sourced replay.
\end{abstract}

\tableofcontents

\section{Concept}\label{concept}

The product is a native multiplayer game client, packaged for Windows,
macOS, Linux on x86\_64 and ARM64. It is consumed by individual end users
who pick up the application, register or log back in, and look for an
opponent. The expected interaction pattern is a few short sessions per
day, each lasting between three and ten minutes; latency matters because
moves must visibly land on the opponent's table within a fraction of a
second to feel responsive, while the underlying card mechanics tolerate
neither lost nor reordered actions.

From a distribution standpoint the project is interesting not because the
business domain demands large-scale parallelism --- a single match is
trivially small --- but because the union of all matches in flight imposes
constraints that no single machine could satisfy at once. Three forces in
particular push the architecture outward:

\begin{itemize}[leftmargin=1.4em]
  \item \textbf{Geography.} A player in Bologna and one in San Francisco
  cannot both share a server that gives them sub-200~ms feedback. The
  system must therefore be \emph{regionally partitioned}, with every
  region holding its own independent compute and database, while still
  presenting a coherent global identity space.
  \item \textbf{Survivability.} Losing a single virtual machine in the
  middle of a match is unacceptable: the players want to keep playing
  rather than restart from scratch. The match state must be \emph{durable}
  on stable storage and \emph{reconstructible} on any healthy node of the
  cluster.
  \item \textbf{Elasticity.} Player traffic is bursty (peaks in the
  evening, dips overnight). The control plane has to scale horizontally;
  no node may sit on session-bound state that would prevent it from being
  killed and replaced.
\end{itemize}

Throughout the document we will refer to two Git repositories: the
\emph{backend} (\texttt{scopa2-backend}, Laravel 12 on Octane and
FrankenPHP) and the \emph{client} (\texttt{scopa2-game}, Godot 4.6 in
C\#/.NET~8). Their
boundaries are intentional: the only contract between client and backend
is a small JSON-over-HTTPS API plus a Pusher-compatible WebSocket
protocol, which makes either side independently replaceable.

\section{Requirements}\label{requirements}

This section captures the functional and non-functional commitments we
made before any code was written, and then projects them onto the
distributed-systems features the course put forward as relevant.

\subsection{Functional Requirements}

\begin{enumerate}[label=F\arabic*., leftmargin=2em]
  \item Players must be able to register and authenticate. The credentials
  must survive across client restarts so that returning users land back in
  the main menu without a fresh login.
  \item Players must be able to enter a matchmaking queue. The system
  pairs two players whose ELO ratings are close enough; the tolerance
  must widen with waiting time so that low-traffic moments still produce
  matches.
  \item Once paired, the two players must play a complete game of
  \emph{Scopa Scientifica} extended with the \emph{Santi} mechanic, which
  lets a player spend captured cards to obtain consumable abilities that
  temporarily mutate cards on the table or in hand.
  \item The server must validate every move against the current game
  state and reject any illegal action. Clients are untrusted.
  \item Every move must reach the opponent in near-real time through a
  pub/sub channel; both players must always observe identical state.
  \item A player whose client crashes mid-game must be able to reconnect
  and resume the same match exactly where it was left.
  \item At the end of the game both players' ELO ratings must be updated
  atomically with the final outcome.
\end{enumerate}

\subsection{Non-Functional Requirements}

\begin{enumerate}[label=NF\arabic*., leftmargin=2em]
  \item \textbf{Latency.} The median time between a client's action POST
  and its broadcast on the partner's WebSocket must remain below 200~ms
  when both players share a region.
  \item \textbf{Single-pod resilience.} The death of any one application
  pod or WebSocket pod must be invisible to the player apart from a brief
  visual stutter; the system must recover in less than five seconds.
  \item \textbf{Single-node database resilience.} The crash of one of the
  three database nodes in a region must not lose any committed write nor
  cause a write to fail.
  \item \textbf{Region survivability.} The catastrophic loss of an entire
  region must leave the surviving region able to serve all global
  traffic. A bounded amount of data may be lost (recovery point
  objective on the order of one second), but no global outage is
  acceptable.
  \item \textbf{Operability.} Releases must be reproducible and require no
  manual intervention from a developer's workstation.
\end{enumerate}

\subsection{Implementation Requirements}

The choice of Laravel as the application framework, Godot with C\# as
the client, and Kubernetes (in the K3s flavour) as the compute substrate
was driven partly by the team's familiarity and partly by the desire to
combine technologies that already implement the lower layers of
distributed coordination correctly. In particular we deliberately picked
YugabyteDB for storage because it embeds a Raft-based replicated log
under a PostgreSQL wire protocol, and Soketi for the realtime tier
because it is wire-compatible with Pusher and ships with a Redis adapter
that handles the cross-pod fan-out problem in a single configuration
flag. 

\subsection{Relevant Distributed-Systems Features}

\begin{description}[leftmargin=1.6em]
  \item[Location transparency.] The client should not need to know which
  region or which pod is serving it; the discovery layer must hide that.
  \item[Replication transparency.] A player's data should look the same
  regardless of which database replica answers the query.
  \item[Failure transparency.] Partial failures of pods, master Redis
  nodes, or even of an entire region should be absorbed by the
  infrastructure with at most a brief visual artefact.
  \item[Openness.] Every protocol used externally is open: HTTPS, the
  Pusher WebSocket protocol, the PostgreSQL wire format. There is no
  proprietary lock-in.
  \item[Scalability.] Horizontal scaling on the application tier; sharded
  storage on the data tier; queue-shape ZSET for matchmaking.
  \item[Concurrency control.] Two pieces of shared state are subject to
  contention: the matchmaking queue (every join writes a new entry) and
  the per-game event log (every move appends a new row). Each is
  protected differently --- the first by atomic Lua, the second by a
  database row lock plus a uniqueness constraint.
  \item[Fault tolerance.] Single-node crashes, network partitions,
  whole-region outages.
  \item[Security.] TLS at the edge, bearer-token authentication,
  authorised private broadcast channels.
\end{description}

\section{Design}\label{design}


\subsection{Architectural Style}\label{architecture}

The architectural choice that matters most in Scopa2 is a recovery one
rather than a stylistic one: every match is modelled as an
\emph{event-sourced deterministic state machine}. The authoritative
state of a game is not a snapshot stored in a database row but the
totally ordered sequence of moves persisted in the
\texttt{game\_events} log. The in-memory \texttt{ScopaEngine} is pure:
given the same seed and the same event list it always produces the
same \texttt{GameState}.

This single decision shapes the rest of the system. Application pods
can be stateless and interchangeable because any of them can
reconstruct any match by replaying its log. Recovery from a crashed
pod or even from a region-wide outage reduces to fetching the log from
a surviving cluster and replaying it; most of the resilience
properties of the system fall out as a corollary of determinism rather
than as the product of dedicated recovery code
(see \Cref{fault-tolerance}).

Two ancillary patterns surround the core. An event-driven
publish/subscribe bus handles the realtime fan-out: validated moves are
emitted as domain events on Pusher-compatible channels and delivered to
whichever pod holds each player's WebSocket. A shared Redis dataspace
holds the matchmaking queue, on which any application pod can operate
without keeping any in-memory state itself. Both choices serve the same
end as the event-sourced engine: keep the application tier stateless so
the system survives the loss of any of its parts.

\subsection{Infrastructure}\label{infrastructure}

\figtikz{topology}{Global symmetric topology. Each region (EU,
North America) hosts an independent stack: an NGINX ingress in front of
a horizontally-scaled fleet of FrankenPHP application pods, a pair of
Soketi WebSocket pods, a three-node Redis cluster with three Sentinels,
and a YugabyteDB cluster of three masters and three tablet servers
configured for intra-region Raft replication and asynchronous
cross-region XCluster replication. Clients reach the closest region via
a latency-aware client-side endpoint selector that pings each region's
\texttt{/up} probe before connecting.}

The deployment substrate is K3s, a lightweight Kubernetes distribution
that we operate on a single physical host with two regional namespaces
(\texttt{region-eu} and \texttt{region-us}) plus a shared
\texttt{db}/\texttt{cache-eu}/\texttt{cache-us} set of namespaces for
the stateful tiers. Every component is templated through a single Helm
chart whose region-specific overrides live in dedicated values files.

\paragraph{Ingress.} A single NGINX ingress per region terminates TLS
and routes traffic: HTTP requests with path \texttt{/api/*} are
forwarded to the application service, while requests on \texttt{/ws/*}
land on the Soketi service. The same ingress also exposes the
\texttt{/up} liveness probe consumed by the client-side endpoint
selector.

\paragraph{Application tier.} The Laravel application is packaged into a
multi-stage Docker image built on \texttt{dunglas/frankenphp}: a single
PHP~8.3 binary running in worker mode keeps the framework resident
across requests, eliminating the per-request bootstrap cost typical of
PHP-FPM deployments. The pod is started by Supervisor, which runs two
processes side-by-side: \texttt{octane:frankenphp} with five workers per
pod and \texttt{queue:work} for background jobs. A horizontal pod
autoscaler ramps the deployment between two and ten replicas based on a
70\% CPU target.

\paragraph{Realtime tier.} Soketi runs as a separate deployment of two
pods per region. The internal traffic from Laravel to Soketi never
leaves the cluster: the application uses Laravel's Pusher broadcast
driver pointed at the in-cluster service DNS name. The two Soketi
replicas share state through a Redis adapter; they would otherwise
deliver every event only to whichever subset of clients happens to be
attached to the same pod.

\paragraph{Cache and coordination tier.} A Redis cluster of three data
nodes and three Sentinels lives in the regional cache namespace. The
Sentinels are responsible for monitoring the master and triggering
automatic failover when it goes down. Both the application (via the
\texttt{phpredis-sentinel} driver) and Soketi (via its native sentinel
support) talk to the Sentinels rather than to a hard-coded master IP.
We therefore inherit the failover logic for free: when the master
crashes, every client of the cluster receives the new master from the
Sentinels and reconnects.

\paragraph{Data tier.} YugabyteDB is a distributed SQL database that
speaks the PostgreSQL wire protocol but stores tablets through a Raft
replication group of three peers; an additional master service stores
the cluster topology and metadata. We deploy it in the \texttt{db}
namespace via the official Helm chart, configured with three masters
and three tablet servers. The application connects through a
comma-separated list of preferred hosts (the regional ``local'' node
first, the other two as fallbacks), which makes PDO transparently
failover at the connection level if the preferred node dies.

\paragraph{Client-side discovery.} A defining choice of our
infrastructure is that we do \emph{not} delegate region selection to a
Geo-DNS provider. The Godot client carries a static list of regional
endpoints and runs an \texttt{EndpointSelector} on startup that probes
each one three times and scores it by
$\text{avg-ping} + 100 \cdot \text{loss}\%$. The lowest-scoring reachable
endpoint becomes the active one. This decision is discussed at length in
\Cref{availability}: it allows us to perform global routing without
owning DNS authority and to react to outages in seconds rather than in
DNS-TTL-bounded minutes.

\paragraph{CI/CD.} GitHub Actions builds the Docker image, pushes it to
the GitHub Container Registry, and then opens an SSH tunnel through
Tailscale to the K3s host where it runs a Helm upgrade for both
regions. The whole pipeline is triggered on every push to the main
branch and the entire deployment cycle (build, push, deploy, verify) is
typically under seven minutes.

\subsection{Domain Modelling}\label{modelling}

The persistent model is intentionally minimal. The
\texttt{users} table holds the immutable identity (id, username,
\texttt{elo} integer). The \texttt{games} table records the existence
of a match: a UUID primary key, the two player ids (the second of which
may be null until the partner joins), a sixteen-character random
\texttt{seed} that fixes the deck-shuffling RNG, and a status enum that
transitions through \texttt{waiting\_for\_players}, \texttt{playing},
and \texttt{finished}. The \texttt{game\_events} table is the
authoritative log: every row is one move, identified by
\texttt{(game\_id, sequence\_number)} with a uniqueness constraint that
ensures no two writes ever land in the same logical slot. The body of
the row carries the actor (\texttt{p1} or \texttt{p2}) and a textual
notation derived from chess PGN.

The transient model is held by the \texttt{ScopaEngine}, a PHP class
that knows how to take the seed of a game, replay the persisted event
list, and yield a \texttt{GameState}. Crucially the engine is
\emph{pure}: it does not touch the database, the broadcaster, or the
clock; given the same input list it always produces the same output
state. This is the property the course refers to as the
\emph{Piecewise Deterministic Assumption}~\citep{elnozahy2002survey},
and we exploit it as the cornerstone of recovery
(\Cref{fault-tolerance}).

The matchmaking model lives outside the relational database, in Redis.
Two keys cooperate: a sorted set \texttt{matchmaking:queue} whose
members are user ids and whose scores are their ELO ratings, and a
hash \texttt{matchmaking:timestamps} that maps each queued user id to
the Unix timestamp at which they entered. Together they form a
\emph{queue with a priority dimension}: the matchmaker can walk the
sorted set in ELO order, while the timestamps inform how patient each
candidate has already been.

\subsection{Interaction Patterns}\label{interaction}

Three patterns describe the totality of inter-component communication
in Scopa2.

\paragraph{Request/response (synchronous).} The client issues HTTPS
requests against a small REST surface:

\begin{itemize}[leftmargin=1.4em]
  \item \texttt{POST /api/auth/register} --- create a new account.
  \item \texttt{POST /api/auth/login} --- authenticate an existing user.
  \item \texttt{POST /api/matchmaking/join} --- enter the matchmaking queue.
  \item \texttt{POST /api/matchmaking/leave} --- leave the queue.
  \item \texttt{GET\phantom{T} /api/matchmaking/status} --- check the
  current queue state.
  \item \texttt{POST /api/games} --- create a new game.
  \item \texttt{POST /api/games/\{id\}/join} --- join an existing game.
  \item \texttt{GET\phantom{T} /api/games/\{id\}} --- fetch the
  canonical state of a game.
  \item \texttt{POST /api/games/\{id\}/action} --- submit a move.
\end{itemize}

All endpoints carry a \texttt{Bearer} token (Sanctum) which makes every
request unambiguously identity-bound. Responses are JSON.

\paragraph{Publish/subscribe (asynchronous).} The realtime layer follows
the Pusher protocol over WebSockets. The Godot client opens a single
WebSocket per session, authenticates the connection's
\texttt{socket\_id} against the backend's \texttt{/broadcasting/auth}
endpoint, and then subscribes to a small set of private channels:

\begin{itemize}[leftmargin=1.4em]
  \item \texttt{\{user\_id\}\_matchmaking\_result} --- the user-private
  channel where the \texttt{match\_found} event is delivered.
  \item \texttt{game\_\{game\_id\}\_player\_\{user\_id\}} --- a
  channel scoped to one player of one match, on which
  \texttt{game\_state\_updated} events arrive with the player's
  individual view.
  \item \texttt{game\_\{game\_id\}} --- a shared channel for both
  players where round-end and game-end events are broadcast.
  \item \texttt{heartbeat} --- a public channel on which the server
  emits a beat every few seconds; the client uses it as a liveness
  indicator (see \Cref{availability}).
\end{itemize}

\figtikz{seq-matchmaking}{Sequence of the discovery and
matchmaking phase. The client first runs the endpoint selector, then
performs authentication, joins the queue, opens a WebSocket,
authenticates the socket against the backend, subscribes to
\texttt{\{user\_id\}\_matchmaking\_result}, and waits. Once two
compatible players are paired, the Lua matchmaker emits the
\texttt{match\_found} event with the new game id; both clients
subscribe to the game's channels and fetch the initial state.}

\figtikz{seq-action}{Sequence of a single in-game action. The
client POSTs the move to the backend; Octane acquires a row lock on
the game; the engine replays the existing event list, applies the
incoming action, persists the resulting \texttt{GameEvent} with the
next sequence number, and broadcasts the new state to both players'
private channels. Soketi receives the broadcast on one of its pods,
publishes it on the Redis adapter, and every other Soketi pod that
holds a connection to either player delivers the event over the
WebSocket.}

\paragraph{Authoritative state transfer.} On every reconnect or
explicit re-sync the client issues \texttt{GET /api/games/\{id\}} and
the server replays the entire event list, producing the canonical
state for the requesting player. The protocol is therefore
\emph{idempotent in the recovery direction}: any client can drop a
message on the wire and recover by asking for the full state.

The serialisation format on the wire is JSON throughout, with a couple
of pragmatic conventions: the opponent's hand is exposed as a list of
``X'' placeholders (so that the message structure is identical for
both players but the cards themselves are hidden), and card mutations
are carried as a flat map from the original card code to its currently
visible code, e.g. \texttt{\{"4C":"4D"\}} after a \emph{San~Biagio}
effect (a Santo that swaps the suit of a chosen card on the table,
here turning a four of clubs into a four of diamonds so the holder of
a four of diamonds can capture it).

\subsection{Behaviour}\label{behaviour}

\figtikz{state-match}{State machine of a match. A game is born in
\texttt{WAITING\_FOR\_PLAYERS}, transitions to \texttt{PLAYING} once
both players have joined, alternates internally through several rounds
(\texttt{deal} $\to$ \texttt{play} $\to$ \texttt{advance round}), and
ends in \texttt{FINISHED} once one of the players reaches the winning
score. \emph{Santi} effects and shop transactions do not advance the
turn; only the play of a card from a player's hand does.}

The engine reconstructs the in-memory state by replaying the
persisted event list from the very first event. The deterministic
shuffle is anchored by a CRC32 of the game's seed for the first round
and of \texttt{seed\_round\_N} for subsequent rounds, so the same
inputs always yield the same dealt cards. The
PGN-inspired notation distinguishes three kinds of action: a card
play (\texttt{7Dx7C\#}, where the optional \texttt{\#} marks a
\emph{scopa}), a Santo invocation (\texttt{@SAN[params]}), and a shop
purchase (\texttt{\$SAN(3C+4D)}). Only card plays advance the turn
counter; Santi and shop interactions modify state without giving up the
turn. When a round ends --- both hands empty and the deck exhausted ---
the engine awards the remaining table cards to the player who last
captured, then computes the round score (denari count, settebello,
primiera, allungo, and scopas). If neither player has reached the
winning score the engine reshuffles and starts a new round; otherwise
it marks the game as finished and triggers the game-end callback.

The engine exposes two callbacks (\texttt{onRoundEnded} and
\texttt{onGameEnded}) which the controller wires to the corresponding
domain events. 

\subsection{Data and Consistency Issues}\label{data-and-consistency-issues}

This is the section where the system's most interesting trade-offs
live. We discuss them through three lenses --- ordering, atomicity, and
the CAP/PACELC envelope --- before commenting on why the more elaborate
machinery of vector clocks and global snapshots is not present.


\paragraph{Atomic matchmaking.} The single place where multiple
application pods race for the same shared resource is the matchmaking
queue. A naive implementation would issue a sequence of Redis commands
from PHP --- inspect the ZSET, decide on a pair, remove both members,
create the game --- but every step is a network round-trip and any
interleaving between two pods can leave the system in a degenerate
state (one player double-matched, or worse, removed but never matched).
We solve the problem by uploading the entire matchmaking decision to
Redis itself as a Lua script. The script walks the queue in ELO order,
computes for each candidate the maximum tolerable ELO gap as a
function of the candidate's wait time, finds a partner that satisfies
both candidates' windows symmetrically, removes both atomically through
\texttt{ZREM} and \texttt{HDEL}, and returns the pair to the caller.
Redis guarantees that the script runs to completion without
interleaving any other command on the database, which makes the
critical section a true atomic step. \Cref{lst:lua-matchmaker} shows
the heart of the algorithm.

\begin{lstlisting}[language={},mathescape=true,caption={Atomic
matchmaking section, shown as pseudocode. The real implementation runs
server-side inside Redis as a single Lua script, which by Redis'
execution model is guaranteed to run to completion without
interleaving any other command --- this is what makes the critical
section a true atomic step regardless of how many application pods
race for the same queue.},label={lst:lua-matchmaker}]
players  = ZSET matchmaking:queue, sorted by ELO
paired   = $\emptyset$
for each p1 in players (in ELO order):
    if p1 $\in$ paired: continue
    w1 = elo_window(wait_time(p1))    // widens with patience

    for each p2 in players coming after p1:
        if p2 $\in$ paired: continue
        w2 = elo_window(wait_time(p2))

        if |p1.elo - p2.elo| $\le$ min(w1, w2):
            ZREM queue p1, p2
            HDEL timestamps p1, p2
            emit pair (p1, p2)
            paired $\leftarrow$ paired $\cup$ {p1, p2}
            break
\end{lstlisting}

\paragraph{CAP and PACELC.} The system spans more than one consistency
domain and treats them differently. Within a region the data tier is
\emph{strongly consistent}: a YugabyteDB tablet leader collects a Raft
majority before acknowledging a write, which means we sit in the CP
corner of the CAP triangle \citep{brewer2000towards, gilbert2002cap}. A
write to a match's event log either commits cluster-wide or fails
visibly to the client; we never accept stale reads or split-brain
writes. Between regions, by contrast, replication is asynchronous
(XCluster): the EU and US clusters reconcile changes in the background,
which means a global reader can briefly observe a US-side game from an
EU vantage point, and recovery from a region-wide outage may lose at
most the tail of an in-flight XCluster batch. We accept this AP
posture because no live match ever spans regions (a game is always
fully owned by one region) and the lost data is bounded to be the last
~1~s of cross-region traffic, which in practice means at most one or
two writes per region. The realtime tier, finally, is consciously
\emph{best-effort}: a broadcast event that does not reach a client is
acceptable as long as the client can pull the canonical state with
\texttt{GET /games/\{id\}}. Pub/sub delivery is therefore AP, the
relational storage is CP, and the boundary is the controller layer
that performs both writes in order.

\citet{abadi2012consistency} adds the PACELC observation that even when
there is no partition the system must still pick between latency and
consistency. Our choice is unambiguous on the intra-region side --- we
pay the few extra milliseconds of Raft quorum to keep the write linearisable.
On the inter-region side we pick latency: we will not block a
European write on the acknowledgment of its American mirror.

\paragraph{Why no Chandy-Lamport snapshots?} \citet{chandy1985distributed}
solve the problem of taking a consistent snapshot of a system in motion
without halting it. We considered the construction during design and
concluded that we did not need it: the only globally consistent cut we
ever need is the per-match event log, which is already a totally
ordered sequence held in a single Raft group. A backup strategy that
naively snapshots each shard independently would risk inconsistent
cross-shard cuts, but YugabyteDB itself provides cluster-wide
backup/restore semantics, so the problem is solved one layer below us.

\subsection{Fault Tolerance}\label{fault-tolerance}

We design for the following failure modes, in roughly decreasing order
of frequency.

\paragraph{Single application pod death.} The application tier is
strictly stateless: every request reconstructs the game from the event
log and tries to acquire a row lock on the game record. Killing a pod
mid-request lets Kubernetes reschedule it; the ingress drops the
failed pod from its rotation and forwards the next attempt to a
healthy peer. The HTTP retry on the client side is the safety net.

\paragraph{Single Soketi pod death.} A client whose Soketi pod dies
loses its WebSocket. The Godot client treats this as a connection
event and reconnects to the same region's WebSocket service; the
ingress routes it to a different Soketi pod. The new pod re-uses the
same Redis adapter, so any event that arrives during the reconnect
window will be delivered to whichever pod the client lands on.

\paragraph{Redis master crash.} The three Sentinels constantly probe
the Redis master. When they observe enough consecutive failures (the
\texttt{down-after-milliseconds} value) they hold an election and
promote one of the replicas; the new master's identity is then
returned to clients through the standard Sentinel discovery protocol.
Both the application (via \texttt{phpredis-sentinel}) and Soketi pick
up the change transparently. A small subset of in-flight Pub/Sub
messages may be lost during the gap; the system tolerates this by
relying on \texttt{GET /games/\{id\}} as the canonical resync path.

\paragraph{YugabyteDB tablet leader crash.} Each tablet of the data
tier is replicated through a Raft group of three peers. The crash of
the leader is detected by the followers when they stop receiving
heartbeats; one of them wins a leader election and resumes serving the
shard. The PostgreSQL driver retries the failed query once on the next
host in its preferred-host list, and the application sees at most a
brief latency spike (typically a few hundred milliseconds).

\paragraph{Whole-region outage.} The catastrophic case --- the entire
EU region disappearing from the network --- is the most expensive to
handle. Three things happen in cascade. First, the client's heartbeat
channel falls silent; after fifteen seconds the
\texttt{TriggerFailover} routine in the Godot client marks the region
as failed and connects to the next reachable endpoint in its priority
list (here, North America). Second, the new region's application
pods have not been serving the player's games, so a re-fetch of the
match returns the latest state replicated across XCluster. Third, the
client subscribes to the game's channels on the new region and
resumes play. The crucial property is that no operator action is
required: the global failover happens client-side.

\paragraph{Recovery as deterministic replay.} The whole design rests on
the property that a game can always be reconstructed from its event
log. Conceptually this is a form of \emph{pessimistic message
logging}~\citep{elnozahy2002survey}: every external input (every move)
is persisted synchronously to stable storage before its effect is
observable to anyone else. The state machine is deterministic
(piecewise deterministic assumption holds because the random number
generator is seeded from the game id), so replaying the persisted
inputs through the engine yields the same state every time. The
trade-off is that the replay cost grows linearly with the length of
the game; in practice a Scopa game is short enough that we never
exceed a few hundred events, so a full replay is cheap. A more
elaborate implementation would periodically snapshot the engine state
to bound the replay cost, but for our match lengths the unbounded
replay is acceptable.

\subsection{Availability}\label{availability}

Availability in Scopa2 is the composition of several independent
mechanisms whose union yields the user-perceived experience of a
continuously running game.

\paragraph{Stateless replication.} The application tier is configured
with a horizontal pod autoscaler whose minimum replica count is two:
even at idle the deployment has more than one survivor. Rolling
updates use \texttt{maxUnavailable: 1} and \texttt{maxSurge: 1}, which
means a deployment under traffic always has a healthy quorum of pods
serving requests.

\paragraph{Heartbeat-driven failover.} The realtime layer publishes a
beat on the public \texttt{heartbeat} channel every few seconds. The
Godot client tracks the time since the last beat received and, after
a fifteen-second silence, triggers a failover. This avoids the
classical problem of trying to distinguish a slow server from a
disconnected one: as long as the heartbeat keeps coming the client
assumes the server is fine even if no game events flow.

\paragraph{Client-side priority list with cooldown.} The
\texttt{EndpointSelector} ranks the regional endpoints once at
startup. After a failover the failed endpoint is added to a set of
``known-bad'' endpoints for the next five seconds, which prevents the
client from flapping between two regions that happen to time out
sequentially. A subsequent failover walks down the priority list.

\paragraph{Channel re-subscription.} On every WebSocket reconnect the
client iterates its set of active channels and re-issues both the
broadcast authentication request and the \texttt{pusher:subscribe}
message. This is what makes the failover seamless: the user's
subscriptions follow them across pods (and across regions).

The result is that an attentive player observes a brief freeze of the
opponent's cards during a failover (typically a second or two) before
the board re-syncs and play resumes. Less attentive players never
notice at all.


\section{Implementation}\label{implementation}

The implementation chapter describes the moving parts of the system in
the order in which a request flows through them, and then concentrates
on the four pieces of code we consider most representative of the
distributed-systems concerns the project addresses: the matchmaking
engine, the game-engine replay, the client's failover logic, and the
configuration that wires the Pusher-compatible bus onto Sentinel-managed
Redis.

\subsection{Technology Stack}

\Cref{tab:stack} groups the chosen technologies by architectural
layer.

\begin{table}[H]
  \centering
  \begin{tabular}{@{}lll@{}}
    \toprule
    \textbf{Layer} & \textbf{Component} & \textbf{Technology} \\
    \midrule
    Client                      & Game UI                & Godot 4.6 (C\# / .NET 8)             \\
    Client networking           & HTTP + WebSocket       & Godot \texttt{HttpRequest}, custom Pusher impl \\
    Edge                        & Ingress                & NGINX (K3s default)                  \\
    Application                 & Web framework          & Laravel 12 on Octane + FrankenPHP    \\
    Application                 & Broadcasting driver    & \texttt{pusher/pusher-php-server}    \\
    Application                 & Redis client           & \texttt{phpredis} + Sentinel adapter \\
    Realtime                    & WebSocket broker       & Soketi (Pusher-compatible)           \\
    Coordination                & Pub/Sub + matchmaking  & Redis 7 + Sentinel                   \\
    Storage                     & Distributed SQL        & YugabyteDB (3 master, 3 tserver)     \\
    Orchestration               & Container runtime      & K3s (lightweight Kubernetes)         \\
    Orchestration               & Packaging              & Helm chart with regional values      \\
    CI/CD                       & Build + deploy         & GitHub Actions, GHCR, Tailscale, SSH \\
    \bottomrule
  \end{tabular}
  \caption{Technology choices, grouped by architectural layer.}
  \label{tab:stack}
\end{table}

\subsection{Repository Layout}

\begin{itemize}[leftmargin=1.4em]
  \item \texttt{scopa2-backend/} --- Laravel application. The relevant
  subtrees are \texttt{app/} (\texttt{Services}, \texttt{GameEngine},
  \texttt{Http/Controllers}, \texttt{Events}, \texttt{Models}),
  \texttt{config/} (Redis Sentinel and broadcasting configuration),
  \texttt{database/migrations/} (the three migrations described in
  \Cref{modelling}), \texttt{helm/} (the chart and per-region values
  files), and \texttt{tests/} (Pest~4 unit and feature suites).
  \item \texttt{scopa2-game/} --- Godot project. The C\# scripts live
  under \texttt{scripts/}; of particular interest are
  \texttt{NetworkManager.cs} (HTTP and WebSocket plumbing,
  failover, channel resubscription), \texttt{EndpointSelector.cs}
  (ping-based region discovery), \texttt{PusherClient.cs} (the
  hand-rolled Pusher protocol implementation), and \texttt{MainGame.cs}
  (UI flow and game-state synchronisation).
\end{itemize}

\subsection{Matchmaking Engine}

The matchmaking subsystem has the two interfaces: an \emph{ingestion}
side (the \texttt{POST /matchmaking/join} handler that pushes an entry
onto the queue and the corresponding leave/status endpoints) and an
\emph{extraction} side (the Lua-driven \texttt{findMatches()} routine
that pairs candidates). The ingestion side is straightforward and
contained in \texttt{MatchmakingController}; the extraction side is
where the interesting distributed logic lives. The Lua script of
\Cref{lst:lua-matchmaker} runs server-side inside Redis, which gives us
two important properties at once: \emph{atomicity} (no other command on
the same Redis can interleave) and \emph{locality} (no
inter-network round-trip per iteration of the inner loop, which matters
because the queue can grow to tens of thousands of entries during peak
traffic).

The ELO window opens dynamically as players wait. Both candidates must
accept each other within their respective windows, which prevents the
pathological case in which a high-ELO player who has just joined would
be paired with a low-ELO player whose window has grown wide enough to
include them but whose own window is still narrow.

\subsection{Game-Engine Replay}

The single most important design property of the backend is that the
authoritative state of any match is a function of the persisted event
log only. \Cref{lst:engine-replay} shows the engine's replay primitive.

\begin{lstlisting}[language=PHP,caption={Deterministic replay loop. The
engine is seeded from the game's seed, then each persisted event is
applied in order; the engine's output is the canonical \texttt{GameState}.},label={lst:engine-replay}]
public function replay(array $historyEvents): GameState
{
    foreach ($historyEvents as $event) {
        $this->applyAction($event->actor_id, $event->pgn_action);
    }
    return $this->state;
}
\end{lstlisting}

Replay is invoked on three occasions: when a client fetches a match
via \texttt{GET /games/\{id\}}, when the action handler reconstructs
the state in order to validate the incoming move, and implicitly during
unit tests. Because the engine is pure, the same event list produces
the same state from any process, in any region, at any time.

\subsection{Action Handling and Concurrency Control}

\Cref{lst:engine-handle} shows the action handler. The transaction
wrapper combined with \texttt{lockForUpdate()} on the \texttt{games}
table acquires a row lock for the duration of the action: any
concurrent move for the same game waits for the lock to be released
before reading the event list. The uniqueness constraint on
\texttt{(game\_id, sequence\_number)} is the safety belt that ensures
the row lock cannot be bypassed even in the presence of a buggy code
path.

\begin{lstlisting}[language=PHP,caption={Action handler. The DB
transaction plus the row lock plus the unique constraint together
provide the \emph{strict serialisation point} that lets us treat
\texttt{sequence\_number} as a Lamport-equivalent total order over a
single match's moves.},label={lst:engine-handle}]
return DB::transaction(function () use ($gameId, $requestData) {
    $game = Game::where('id', $gameId)->lockForUpdate()->firstOrFail();
    $events = GameEvent::where('game_id', $gameId)
        ->orderBy('sequence_number')->get();

    $engine = new ScopaEngine($game->seed, $onRound, $onGame);
    $engine->replay($events->all());

    $engine->applyAction($this->getLoggedPlayerIndex($game), $action);

    if ($engine->getState()->isGameOver) {
        $game->status = GameStateEnum::FINISHED;
        $game->save();
    } else {
        broadcast(new GameStateUpdated($game->id,
            $engine->getState()->toPublicView('p1'), $game->player_1_id));
        broadcast(new GameStateUpdated($game->id,
            $engine->getState()->toPublicView('p2'), $game->player_2_id));
    }

    GameEvent::create([
        'game_id'         => $gameId,
        'sequence_number' => $events->count() + 1,
        'actor_id'        => $this->getLoggedPlayerIndex($game),
        'pgn_action'      => $action,
    ]);
    return response()->json(['status' => 'success']);
});
\end{lstlisting}

\subsection{Realtime Wiring through Sentinel}

The Soketi deployment receives, through environment variables, a JSON
array of Sentinel endpoints rather than a static master IP. The same
list is used both for the adapter (the cross-pod fan-out bus) and for
the application/presence database. The configuration is templated
through Helm so that the EU and US deployments inject their own
regional Sentinels.

\begin{lstlisting}[language=bash,caption={Excerpt of the Soketi
container's environment. The same Sentinel list serves two distinct
roles: as the source of truth for the adapter (cross-pod event
distribution) and for the database (presence/app state).},label={lst:soketi-env}]
- name: SOKETI_ADAPTER_DRIVER
  value: "redis"
- name: SOKETI_ADAPTER_REDIS_INSTANCE_NAME
  value: "mymaster"
- name: SOKETI_ADAPTER_REDIS_SENTINELS
  value: '[{"host":"redis-sentinel-0...","port":26379},
           {"host":"redis-sentinel-1...","port":26379}]'
- name: SOKETI_DB_DRIVER
  value: "redis"
- name: SOKETI_DB_REDIS_SENTINELS
  value: '[{"host":"redis-sentinel-0...","port":26379},
           {"host":"redis-sentinel-1...","port":26379}]'
\end{lstlisting}

The clients of this configuration do not know the master's IP at any
point; they discover it through the Sentinels at connection time and
re-discover it on every failover.

\subsection{Client-Side Failover}

The client's \texttt{NetworkManager} responsibility
 is to keep the user in a connected state without
exposing the underlying complexity. \Cref{lst:endpoint} shows the scoring
formula used by the endpoint selector and the failover policy.

\begin{lstlisting}[language={[Sharp]C},caption={Endpoint scoring and
selection. The penalty term $100 \cdot \text{loss}$ heavily punishes
unreliable endpoints even when their average latency is good; an
unreachable endpoint is given $\texttt{double.MaxValue}$ so it is
always sorted last.},label={lst:endpoint}]
public double Score => IsReachable
    ? AveragePingMs + (PacketLoss * 100)
    : double.MaxValue;

public async Task<ServerEndpoint> SelectBestEndpoint(
        ServerEndpoint[] endpoints) {
    var metrics = await GetEndpointMetrics(endpoints);
    var reachable = metrics.Where(m => m.IsReachable).ToArray();
    if (reachable.Length == 0) return endpoints[0]; // last-resort fallback
    return reachable.OrderBy(m => m.Score).First().Endpoint;
}
\end{lstlisting}

When a failover is triggered (either by exhaustion of HTTP retries or
by absence of heartbeat for fifteen seconds) the client adds the
current endpoint to a ``known-bad'' set, picks the next reachable
endpoint from the priority list, reconnects the WebSocket, re-issues
the broadcast authentication for every active channel, re-subscribes
to those channels, and finally calls \texttt{GET /games/\{id\}} to
resync the board.

\subsection{Continuous Integration and Deployment}

The CI/CD pipeline lives in
\texttt{.github/workflows/deploy.yml}. Every push to \texttt{main}
triggers a Docker build with the commit SHA as the image tag, pushes
the image to the GitHub Container Registry, brings up an ephemeral
Tailscale node, SSHes into the K3s host, and runs
\texttt{helm upgrade} for both \texttt{region-eu} and
\texttt{region-us}. The rollout uses Kubernetes' rolling-update
strategy, so the production deployment never observes a zero-replica
window. The whole pipeline typically takes between five and seven
minutes.

\section{Validation}\label{validation}

The validation strategy is two-pronged. Automated tests provide the
fast feedback loop on a developer's machine; manual fault-injection
tests measure how the deployed system reacts to the failure modes the
architecture is meant to handle.

\subsection{Automated Testing}

The backend test suite is written in Pest~4 and split between unit
tests (which exercise the pure game engine and its support classes)
and feature tests (which spin up the Laravel application against a
SQLite test database and a Redis instance and drive it through HTTP).
Notable coverage:

\begin{itemize}[leftmargin=1.4em]
  \item \texttt{GameEngine/} --- \texttt{ScopaEngineTest},
  \texttt{ScoreCalculatorTest},
  \texttt{ScopaNotationParserTest},
  \texttt{MoveValidatorTest},
  \texttt{PlayerStateTest}, \texttt{GameStateTest},
  plus per-Santo tests (\texttt{SanBiagio}, \texttt{SanPantaleone},
  \texttt{SantaCaterina}).
  \item \texttt{Feature/MatchmakingTest} --- enqueue, dequeue,
  \texttt{isQueued}, the HTTP endpoints, queue-size assertions, and
  duplicate-prevention.
  \item \texttt{Feature/GameTest} --- game creation, join,
  state retrieval (with hand visibility checks),
  card discard, capture, and the boundary conditions for the
  \texttt{action} handler.
  \item \texttt{Feature/AuthTest} --- registration, login, and
  logout flows.
\end{itemize}


\subsection{Acceptance and Chaos Testing}

The deployment was exercised by manual fault injection from outside
the cluster. \Cref{tab:chaos} summarises the test matrix derived from
the project backlog notes. Tests in the realtime and application
tiers passed cleanly; one test in the cache tier exposed a real
problem with our Sentinel quorum configuration, listed in
\Cref{future-works}.

\begin{table}[h]
  \centering
  \begin{tabular}{@{}lp{4cm}p{5cm}l@{}}
    \toprule
    \textbf{Tier} & \textbf{Action} & \textbf{Expected outcome} & \textbf{Result} \\
    \midrule
    Application & Kill one app pod & LB drops failed pod, retry on peer pod returns 200 & \texttt{PASS} \\
    Application & Kill all app pods in a region & Client times out, fails over to other region, state preserved & \texttt{PASS} \\
    Realtime    & Scale Soketi to one replica then back to two & Clients reconnect on the surviving pod & \texttt{PASS} \\
    Cache       & Kill a Redis follower & No visible effect & \texttt{PASS} \\
    Cache       & Kill Redis master & Sentinels elect a new master, traffic resumes & \texttt{FAIL} (see \Cref{future-works}) \\
    Storage     & Kill a TServer leader & Raft re-election, query latency spike $\sim$3~s & not yet automated \\
    Storage     & Reduce TServer to one replica & Writes fail; reads via follower-reads, manual intervention required & not yet automated \\
    Global      & Delete an entire region namespace & Client failover to surviving region, RPO bounded & not yet automated \\
    \bottomrule
  \end{tabular}
  \caption{Chaos test matrix. \texttt{PASS} entries were verified by
  hand by deleting Kubernetes pods or scaling deployments while a
  client was actively playing.}
  \label{tab:chaos}
\end{table}

\section{Deployment}\label{deployment}

The production deployment is a single K3s cluster running on a
dedicated host on the home lab network, reachable over Tailscale by
both the CI pipeline and the development team. The public surface
consists of two HTTPS endpoints, one per region:

\begin{itemize}[leftmargin=1.4em]
  \item \url{https://scopa-eu.murkrowdev.org}
  \item \url{https://scopa-us.murkrowdev.org}
\end{itemize}

The endpoints share the same TLS certificate authority and the same
deployment configuration, modulo the regional values file that
overrides the database's preferred host, the Redis Sentinel host, and
the Soketi service name. \Cref{lst:helm-eu} shows the EU values file
in its entirety; the US file differs only in the namespace prefixes
of the same set of hostnames.

\begin{lstlisting}[caption={Helm values overrides for
the EU region. Every other configuration parameter is inherited from
the shared \texttt{values.yaml}.},label={lst:helm-eu}]
ingress:
  host: "eu.scopa.local"

redis:
  host: "eu-scopa-redis.cache-eu.svc.cluster.local"

soketi:
  host: "scopa2-eu-soketi-service"

env:
  REDIS_SENTINEL_HOST: "eu-scopa-redis.cache-eu.svc.cluster.local"

db:
  localNode: "yb-tserver-0.yb-tservers.db.svc.cluster.local"
  backups:   "yb-tserver-1.yb-tservers.db.svc.cluster.local,yb-tserver-2.yb-tservers.db.svc.cluster.local"
\end{lstlisting}

\figtikz{deployment}{Per-region Kubernetes resource graph. The
ingress fronts a Service that load-balances across the deployment of
FrankenPHP pods; the Soketi deployment exposes a separate Service for
the WebSocket traffic. The HPA controls replica count for the
application deployment. Redis runs as a StatefulSet with three pods
(masters/replicas) and three Sentinel pods. YugabyteDB lives in its
own \texttt{db} namespace as two StatefulSets (masters and tservers)
shared between both regional application namespaces.}

For local development we run a single-node Minikube cluster with
identical Helm charts. Routing between the local host machine and the
in-cluster ingress is provided by Caddy on the developer's machine,
which reverse-proxies \texttt{:8080} and \texttt{:8081} to the
Minikube IP with the appropriate \texttt{Host} header (so the same
NGINX ingress can dispatch to either region just by inspecting the
header). The local dev environment can be reached by the rest of the
team through Tailscale, which exposes the developer's machine on the
shared mesh network without opening any port to the public Internet.

\section{User Guide}\label{user-guide}

The expected workflow for an end user is the following.

\begin{enumerate}[leftmargin=1.6em]
  \item Download the platform-appropriate build of the Scopa2 client
  from the project's releases page. Builds are produced for Windows
  x86\_64, macOS Universal (Intel and Apple Silicon), Linux x86\_64,
  and Linux ARM64.
  \item Launch the client. On first run it will probe the regional
  endpoints, pick the lowest-latency reachable region, and present a
  registration screen.
  \item Enter a username and submit the registration form. On success
  the credentials are persisted locally in the
  user-private \texttt{user://auth.cfg} file and the main menu is
  shown.
  \item Click \emph{Find Match} to enter the ELO-based queue. The
  status bar at the bottom of the screen indicates the chosen region
  and the connection status.
  \item Once a partner is found the client transitions automatically
  to the game board. The first card is drawn within at most a few
  seconds; after that, every move propagates between the two clients
  in well under one second under normal network conditions.
  \item In case of network instability the client will display a
  brief \emph{Server Switch} overlay while it migrates to the
  fallback region. The game resumes automatically once the new
  connection is established.
\end{enumerate}

Logout is available from the main menu and clears the persisted
credentials. There is no password recovery flow at the time of
writing because the current authentication model is username-only;
the password-based flow is one of the items in
\Cref{future-works}.

\section{Self-Evaluation}\label{self-evaluation}

\subsection{Work Split}

The two of us met in person almost every day and worked from the same
workstation for the vast majority of the project. As a consequence the
planning, the architectural decisions, and the bulk of the
implementation --- both on the backend and on the Godot client --- were
carried out jointly: there is no meaningful sense in which a specific
controller or scene was ``owned'' by one of us, and the report itself
was written together.

The split surfaces only in the two areas where the work required
hardware or context that only one of us had at hand:

\begin{itemize}[leftmargin=1.4em]
  \item \textbf{Marco} owned the production deployment side. With
  access to a dedicated home-lab server he set up the K3s cluster,
  the per-region Helm releases, the reverse-proxy and domain
  configuration, and the Tailscale-bridged CI/CD pipeline that
  performs the rolling rollout from GitHub Actions into the cluster.
  \item \textbf{Valerio} owned the client-side recovery manager
  inside the Godot project: the ping-based endpoint selection, the
  heartbeat-driven failover logic, the channel resubscription on
  reconnect, and the manual chaos validation that the resulting flow
  actually masks regional outages from the end user.
\end{itemize>

\subsection{Reflections}

Looking at the project from the perspective of the course's
conceptual grid, the most rewarding finding has been that several
canonical distributed-systems mechanisms do not need to be implemented
\emph{by us} to be \emph{used} effectively. The Raft consensus
algorithm \citep{ongaro2014search} lives inside YugabyteDB, the
master-election protocol is delegated to Redis Sentinel, and the
horizontal scaling and rolling-update behaviour comes from Kubernetes.
What the application owns is the boundary code that glues these
primitives together --- the schema that produces a Lamport-equivalent
ordering, the Lua script that turns matchmaking into an atomic step,
the deterministic engine that lets the system recover by replay rather
than by snapshotting.

The most uncomfortable finding has been that distributed-systems
design can easily produce \emph{honest-looking but wrong} solutions.
The clearest example in our codebase is that the Lua matchmaker
correctly serialises the critical section, yet the script itself is
currently invoked piggy-backed on join requests rather than from a
scheduled background process: during a quiet HTTP minute the queue
can keep growing while no pod is actually trying to pair the waiting
players. The correctness of the primitive does not, on its own,
guarantee the correctness of the system that uses it.

Two large portions of the course syllabus deliberately fall outside
our scope. We do not implement smart contracts, distributed ledgers,
or related zero-trust primitives: the domain is a casual card game in
which the server is the trust anchor and the cost of a Byzantine model
would dwarf the value it adds. Likewise we make no use of code
mobility in the sense of \citet{fuggetta1998understanding}: the
client's behaviour is statically compiled, the server's logic ships in
a Docker image, and we have no use case for shipping executable code
across the boundary. We acknowledge these omissions explicitly rather
than glossing them over.

\section{Future Works}\label{future-works}

The items below collect the distributed-systems extensions we
consider most worthwhile to pursue beyond the current scope of the
project.

\begin{itemize}[leftmargin=1.4em]
  \item \textbf{Periodic checkpointing.} For very long games the
  replay cost grows linearly with the event count. A periodic engine
  snapshot would bound the replay window and allow garbage collection
  of obsolete event log entries.
  \item \textbf{Explicit Lamport stamps.} Even though the database row
  lock provides serialisation, attaching an explicit Lamport
  timestamp to each event would make causal reasoning portable across
  shards if we ever decide to scale a single match across regions.
  \item \textbf{Vector clocks for the matchmaking queue.} The queue
  is currently not replicated across regions; a global pairing
  algorithm would need a causality scheme that handles concurrent
  joins from the two regions safely.
  \item \textbf{Redis Sentinel quorum configuration.} Chaos test 2.2
  exposed a real failure mode: when the master is killed under our
  current configuration the Sentinels do not elect a new master. The
  \texttt{quorum}, \texttt{down-after-milliseconds}, and
  \texttt{parallel-syncs} parameters must be tuned and verified.
  \item \textbf{XCluster failover automation.} A region-wide failure
  is currently handled on the client side but the data tier still
  requires manual promotion of the surviving cluster as the new
  XCluster source of truth.
  \item \textbf{Service-mesh mTLS.} Internal traffic between pods
  flows over the plain K3s overlay. Adding Linkerd would provide
  mTLS automatically without application changes.
\end{itemize}

\bibliographystyle{plain}
\bibliography{references}

\end{document}
